\subsection{Local objects}
\label{subsec:2dloc-local}

Fix an $(\infty,2)$-category $\D$ and classes $I$, $P$ of $1$-morphisms of $\D$. $\D$ may be large (e.g.\ $\D = \Cat$), with small hom-categories; functor $(\infty,2)$-categories out of $\D$ are then formed in a larger universe. No other assumption on $\D$ is made in this subsection.
For $(s\colon x \to y) \in I \cup P$ and $d \in \D$ let $s^{*} := \Hom_{\D}(s,d)\colon \Hom_{\D}(y,d) \to \Hom_{\D}(x,d)$, and for $f\colon d \to d'$ let $f_{*}$ be postcomposition.

\begin{defi}
\label{def:IP-local}
An object $d \in \D$ is \emph{$(I,P)$-local} if every $i^{*}$, $i \in I$, admits a fully faithful right adjoint $i_{*}$, and every $p^{*}$, $p \in P$, admits a fully faithful left adjoint $p_{!}$. So $i^{*}$ is a localization and $p^{*}$ a colocalization of $\infty$-categories.
A $1$-morphism $f\colon d \to d'$ between $(I,P)$-local objects is \emph{$(I,P)$-local} if for every $i \in I$ and $p \in P$ the lax squares
\[
\begin{tikzcd}[column sep=large]
\Hom_{\D}(y,d) \ar[r, "f_{*}"] \ar[d, "i^{*}"'] & \Hom_{\D}(y,d') \ar[d, "i^{*}"] \ar[dl, Rightarrow, shorten=4mm, "="'] \\
\Hom_{\D}(x,d) \ar[r, "f_{*}"'] & \Hom_{\D}(x,d')
\end{tikzcd}
\qquad
\begin{tikzcd}[column sep=large]
\Hom_{\D}(y,d) \ar[r, "p^{*}"] \ar[d, "f_{*}"'] & \Hom_{\D}(x,d) \ar[d, "f_{*}"] \ar[dl, Rightarrow, shorten=4mm, "="'] \\
\Hom_{\D}(y,d') \ar[r, "p^{*}"'] & \Hom_{\D}(x,d')
\end{tikzcd}
\]
(given by associativity $i^{*}f_{*} \simeq f_{*}i^{*}$, resp.\ $f_{*}p^{*} \simeq p^{*}f_{*}$) are right, resp.\ left, Beck--Chevalley, i.e.\ $f_{*}i_{*} \Rightarrow i_{*}f_{*}$ and $p_{!}f_{*} \Rightarrow f_{*}p_{!}$ are invertible.
Let $\iota\colon \Pc_{I,P} \hookrightarrow \D$ be the locally full sub-$(\infty,2)$-category on the $(I,P)$-local objects and $1$-morphisms. When $\D$ needs to be named we write $\Pc_{I,P}(\D)$.
\end{defi}

$\Pc_{I,P}$ is well defined: identities are $(I,P)$-local, and $(I,P)$-local $1$-morphisms are closed under composition and equivalence since mates are compatible with pasting (\cref{lem:mate-correspondence}).
By definition $\Pc_{I,P} = \Pc_{I,\emptyset} \cap \Pc_{\emptyset,P}$. No condition relates $i_{*}$ and $p_{!}$, also when $I \cap P \neq \emptyset$.

By the enriched Yoneda lemma (I3), $d \mapsto \Hom_{\D}(-,d)$ and $s \mapsto \Hom_{\D}(s,-)$ are functors; for $s\colon x \to y$ the latter gives $s^{*}\colon \Hom_{\D}(y,-) \Rightarrow \Hom_{\D}(x,-)$ in $\func(\D,\Cat)$, that is, a functor
\[
R_{s}\colon \D \longrightarrow \func(\Delta^{1},\Cat) \ , \qquad R_{s}(d) = s^{*} \ ,
\]
where $R_{s}(f)$, for $f\colon d \to d'$, is the commutative square with sides $s^{*}$, $s^{*}$, $f_{*}$, $f_{*}$ filled by the associativity equivalence $s^{*}f_{*} \simeq f_{*}s^{*}$ (both sides are $f \circ - \circ s$).

\begin{thm}
\label{thm:local-pullback}
The square
\begin{equation}
\label{eq:pullback}
\begin{tikzcd}[column sep=large]
\Pc_{I,P} \ar[r] \ar[d, "\iota"'] & \prod_{I}\func(\Refl,\Cat) \times \prod_{P}\func(\Corefl,\Cat) \ar[d, "\prod_{I}\ell^{*} \times \prod_{P}\rho^{*}"] \\
\D \ar[r, "(R_{s})_{s}"'] & \prod_{I \sqcup P}\func(\Delta^{1},\Cat)
\end{tikzcd}
\end{equation}
is a pullback of $(\infty,2)$-categories, for a unique top map. Its component at $i \in I$ sends $d$ to the reflection $i^{*} \dashv i_{*}$, and its component at $p \in P$ sends $d$ to the coreflection $p_{!} \dashv p^{*}$.
\end{thm}

So $d$ is $(I,P)$-local if and only if the representable $\Hom_{\D}(-,d)$ takes $I$ to left adjoints of reflections and $P$ to right adjoints of coreflections, and similarly for $1$-morphisms, with the Beck--Chevalley conditions of \cref{lem:mates}.
The exchange of $\ell$ and $\rho$ reflects the contravariance of $\Hom_{\D}(-,d)$: making $i$ reflective asks $i^{*}$ to be the left adjoint of a reflection, and making $p$ coreflective asks $p^{*}$ to be the right adjoint of a coreflection.

\begin{proof}
By \cref{lem:mates-refl} for $\E = \Cat$, the right-hand map is a locally full sub-$(\infty,2)$-category, hence so is its pullback along $(R_{s})_{s}$ (\cref{rmk:lax-epi}). It is the locally full sub-$(\infty,2)$-category of $\D$ on those $d$, resp.\ $f\colon d \to d'$, for which every $R_{i}(d)$, resp.\ $R_{i}(f)$, lies in the image of $\ell^{*}$ and every $R_{p}(d)$, resp.\ $R_{p}(f)$, in the image of $\rho^{*}$. We compute these images; as a locally full sub-$(\infty,2)$-category is determined by its objects and $1$-morphisms (\cref{subsec:prelim-conventions}), this proves the claim.

By \cref{lem:mates-refl}, $R_{i}(d) = i^{*}$ lies in the image of $\ell^{*}$ if and only if $i^{*}$ has a right adjoint $i_{*}$ with invertible counit, i.e.\ a fully faithful one, and $R_{p}(d) = p^{*}$ in that of $\rho^{*}$ if and only if $p^{*}$ has a fully faithful left adjoint $p_{!}$. So the objects in the image are the $(I,P)$-local ones.
Let $f\colon d \to d'$ with $d$, $d'$ local. $R_{i}(f)$ is a $1$-morphism from $u = i^{*}$ to $u' = i^{*}$ in $\func(\Delta^{1},\Cat)$ with $f_{0} = f_{*}$ on $\Hom_{\D}(y,-)$ and $f_{1} = f_{*}$ on $\Hom_{\D}(x,-)$. By \cref{lem:mates-refl}(2), \cref{lem:mates}(1) applies and reads it as the lax square with $u$ vertical and $\varphi\colon u'f_{0} \simeq f_{1}u$, here the associativity $i^{*}f_{*} \simeq f_{*}i^{*}$: this is the first square of \cref{def:IP-local}, and $R_{i}(f)$ lies in the image of $\ell^{*}$ if and only if its right mate $f_{*}i_{*} \Rightarrow i_{*}f_{*}$ is invertible. Dually, \cref{lem:mates}(2) reads $R_{p}(f)$ as the lax square with $u = p^{*}$ horizontal and $\psi\colon f_{1}u \simeq u'f_{0}$, here $f_{*}p^{*} \simeq p^{*}f_{*}$: the second square of \cref{def:IP-local}, which lies in the image of $\rho^{*}$ if and only if its left mate $p_{!}f_{*} \Rightarrow f_{*}p_{!}$ is invertible.
No choice is involved: a commutative square has one underlying lax square with $u$ on a given side, and the $2$-morphism is the one of $R_{s}(f)$. (The Beck--Chevalley condition depends on this $2$-morphism, not only on the edges.)
\end{proof}

The reason for the Beck--Chevalley conditions: a lift of $R_{i}(f)$ along $\ell^{*}$ is a $1$-morphism of $\func(\Refl,\Cat)$, i.e.\ a transformation from the reflection $i^{*} \dashv i_{*}$ at $d$ to the one at $d'$ with components $f_{*}$. Besides its naturality square $R_{i}(f)$ at $i^{*}$, it has an \emph{invertible} naturality square at $i_{*}$, and in $\htwo\Cat$ this square is the right mate of $R_{i}(f)$ (\cref{rmk:mates-2cat}). Dually for $p$, with $p_{!}$ and the left mate.

The case $\D = \Pc_{\emptyset,\{s\}}$, i.e.\ every object and $1$-morphism is $\{s\}$-local, characterizes coreflective $1$-morphisms representably:

\begin{lem}
\label{lem:repr-refl}
A $1$-morphism $s$ of $\D$ is coreflective if and only if $\Pc_{\emptyset,\{s\}} = \D$, and reflective if and only if $\Pc_{\{s\},\emptyset} = \D$.
\end{lem}

\begin{proof}
The enriched Yoneda embedding $\D\opcat \to \func(\D,\Cat)$, $d \mapsto \Hom_{\D}(d,-)$, is fully faithful (I3). It takes an adjunction $l \dashv r$ of $\D$ to $r^{*} \dashv l^{*}$ with the same unit and counit, and every adjunction between representables arises this way. So $s$ is coreflective if and only if $s^{*}$ is a right adjoint with invertible unit in $\func(\D,\Cat)$.
By \cref{thm:adjoints-in-functor-cats}, $s^{*}$ is a right adjoint if and only if every component $s^{*}\colon \Hom_{\D}(y,d) \to \Hom_{\D}(x,d)$ is one and every transposed naturality square is left Beck--Chevalley; these squares are the second squares of \cref{def:IP-local}. A $2$-morphism of $\func(\D,\Cat)$ is invertible if and only if its components are, so the unit is invertible if and only if every $s^{*}$ has a fully faithful left adjoint. The reflective case is dual.
\end{proof}

\begin{defi}
\label{def:enriched-adjoint}
A functor $G\colon \C \to \D$ of $(\infty,2)$-categories \emph{has a left adjoint} if there are a functor $F\colon \D \to \C$ and a transformation $\eta\colon \mathrm{id}_{\D} \Rightarrow GF$ such that
\[
\Hom_{\C}(Fy,q) \xrightarrow{\ G\ } \Hom_{\D}(GFy,Gq) \xrightarrow{\ \eta_{y}^{*}\ } \Hom_{\D}(y,Gq)
\]
is an equivalence for all $y \in \D$ and $q \in \C$. These equivalences are natural in $y$ and $q$, as functors $\D\opcat \times \C \to \Cat$.
\end{defi}

In the presentable case (I4) provides such left adjoints. In general $\iota\colon \Pc_{I,P} \to \D$ need not have one; when it does, the adjoined adjoints are precomposition:

\begin{prop}
\label{prop:local-adjoint-refl}
Suppose $\iota\colon \Pc_{I,P} \to \D$ has a left adjoint $L$. Then $Li$ is reflective for $i \in I$ and $Lp$ is coreflective for $p \in P$, and $i_{*} \simeq ((Li)^{L})^{*}$ and $p_{!} \simeq ((Lp)^{R})^{*}$ on $\Hom_{\D}(-,\iota q) \simeq \Hom_{\Pc_{I,P}}(L-,q)$.
\end{prop}

\begin{proof}
The equivalences $\Hom_{\Pc_{I,P}}(Ly,q) \simeq \Hom_{\D}(y,\iota q)$ are natural in $y$ and $q$. They identify $(Li)^{*}$ at $q$ with $i^{*}$ at $\iota q$, and postcomposition by $g\colon q \to q'$ with $(\iota g)_{*}$. So every object and $1$-morphism of $\Pc_{I,P}$ is $\{Li\}$-local, and $Li$ is reflective by \cref{lem:repr-refl}; likewise for $Lp$. Under the Yoneda embedding the reflection $(Li)^{L} \dashv Li$ goes to $(Li)^{*} \dashv ((Li)^{L})^{*}$, so $((Li)^{L})^{*}$ is the right adjoint $i_{*}$ of $i^{*}$; likewise for $p$.
\end{proof}

By \cref{rmk:mates-2cat}, for $g\colon q \to q'$ in $\Pc_{I,P}$ the squares expressing that $g_{*}$ commutes with precomposition by $(Li)^{L}$ and $(Lp)^{R}$ are the mates of the squares of \cref{def:IP-local}: the Beck--Chevalley conditions say that $\iota g$ commutes with precomposition by the adjoined adjoints.

Making $I$ reflective and $P$ coreflective commutes, also on the level of local objects, as soon as the relevant left adjoint exists.

\begin{prop}
\label{prop:2dloc-commute}
Suppose $\iota_{I}\colon \Pc_{I,\emptyset} \to \D$ has a left adjoint $L_{I}$. Then $\iota_{I}$ restricts to an equivalence from $\Pc_{\emptyset,L_{I}(P)}(\Pc_{I,\emptyset})$ onto $\Pc_{I,P}$. Dually, if $\iota_{P}\colon \Pc_{\emptyset,P} \to \D$ has a left adjoint $L_{P}$, then $\iota_{P}$ restricts to an equivalence from $\Pc_{L_{P}(I),\emptyset}(\Pc_{\emptyset,P})$ onto $\Pc_{I,P}$.
\end{prop}

\begin{proof}
As in the proof of \cref{prop:local-adjoint-refl}, $\Hom_{\Pc_{I,\emptyset}}(L_{I}y,q) \simeq \Hom_{\D}(y,\iota_{I}q)$ identifies $(L_{I}p)^{*}$ at $q$ with $p^{*}$ at $\iota_{I}q$, and $g_{*}$ with $(\iota_{I}g)_{*}$. So $q$ is $(\emptyset,L_{I}(P))$-local if and only if $\iota_{I}q$ is $(\emptyset,P)$-local, and likewise for $1$-morphisms; the image is therefore $\Pc_{I,\emptyset} \cap \Pc_{\emptyset,P} = \Pc_{I,P}$. The other case is dual.
\end{proof}

\begin{exmp}
\label{ex:local-cat}
Let $\D = \Cat$ and $P = \{u\}$ for $u\colon X \to Y$. Then $\Pc_{\emptyset,\{u\}}$ is the locally full sub-$(\infty,2)$-category of $\Cat$ on the $\C$ for which $u^{*}\colon \func(Y,\C) \to \func(X,\C)$ has a fully faithful left adjoint $u_{!}$, and the $F\colon \C \to \C'$ for which the commutative square
\[
\begin{tikzcd}
\func(Y,\C) \ar[r, "u^{*}"] \ar[d, "F_{*}"'] & \func(X,\C) \ar[d, "F_{*}"] \ar[dl, Rightarrow, shorten=3mm, "="'] \\
\func(Y,\C') \ar[r, "u^{*}"'] & \func(X,\C')
\end{tikzcd}
\]
is left Beck--Chevalley, i.e.\ $u_{!}F_{*} \Rightarrow F_{*}u_{!}$ is invertible.
If $u$ is fully faithful, a pointwise left Kan extension along $u$ is fully faithful, so every $\C$ with colimits indexed by the comma $\infty$-categories $X \times_{Y} Y_{/y}$ lies in $\Pc_{\emptyset,\{u\}}$. But $u_{!}$ need not be pointwise: for $u\colon \{1\} \to \Delta^{1}$, a space $\C$ with two points lies in $\Pc_{\emptyset,\{u\}}$, as $u^{*}$ is an equivalence, but pointwise left Kan extension along $u$ would require an initial object.
Dually, $\Pc_{\{u\},\emptyset}$ consists of the $\C$ for which $u^{*}$ has a fully faithful right adjoint $u_{*}$, and the $F$ with $F_{*}u_{*} \Rightarrow u_{*}F_{*}$ invertible.
By \cref{thm:2dloc-local} these are $\Cat\corefl{u}$ and $\Cat\refl{u}$.
\end{exmp}
